\section{Uncertainties and validation}
\label{sec:validation}

A closure test measures how well a specified deformation of the known simulated truth is recovered, with the
prior, the initial truth distribution, held fixed.
Coverage, how often intervals contain their target over repeated experiments, tests the estimator and the interval
construction together, and the coverage tested here is pointwise, for each reported quantity separately, not the
stronger simultaneous coverage of all of them.
The closure, deformation, and coverage studies reweight the true signal of the simulated events used as
pseudo-data, which changes their reconstructed and true distributions together and keeps the detector response at
fixed truth.
Throughout, the prior used to unfold them stays at Tune~v1.
The envelope of Sec.~\ref{sec:uncertainty-status} instead reweights the prior used to unfold the data.
Every study uses the analysis simulation's detector response.

\subsection{The reported covariance}
\label{sec:cov16}

The reported covariance of the sixteen-cell result is $C_{\mathrm{total}}+C_{\mathrm{resp}}$.
Unlike the envelope of Sec.~\ref{sec:uncertainty-status}, its model term is a covariance built from three generator
shifts.
The term $C_{\mathrm{total}}$ sums the statistical, flux, interaction, detector, lateral and normalization
uncertainties of the joint-test estimator and $C_{\mathrm{model}}$.
The model term $C_{\mathrm{model}}$ is built from the change in the unfolded data when the prior is reweighted to
GENIE + MEC, NuWro and GiBUU in turn; it represents unscaled generator-shape dependence, not a calibrated bound, and
its coverage at truths held out of its construction has not been validated.
The response term $C_{\mathrm{resp}}$ is the response of the result to a one-standard-deviation scale of the
reconstructed recoil energy in the simulation, for protons, mesons, electromagnetic and other deposits, but it does not cover neutron or muon deposits, nor
species, energy-dependence or resolution effects beyond a scale, and it enters only this result, its fake-data tests
and its comparison in Sec.~\ref{sec:gen16}.
The median per-cell relative uncertainty is \fvval{sigmed-C_total_resp}\%, and \fvval{sigmed-C_total}\% without
$C_{\mathrm{resp}}$.
The note gives the cells, the central values, the per-group uncertainties and the correlation matrix of the
reported covariance.

\subsection{Two-dimensional benchmark}
\label{sec:repro2d}
\label{sec:repro2d-pub}

The two-dimensional result is an end-to-end test of the unfolding against a measurement from the same data.
Its integrated cross section is \SI{\sigTwoD}{cm^2/nucleon}, compared with \SI{\sigTwoDpaper}{cm^2/nucleon} in the
published analysis, a ratio of $\ratioTot$.
The standardized bin differences formed with the published uncertainties have an RMS of $\pullRMS$.
Against the full published covariance, the descriptive covariance-weighted distance is
$\chi^{2}/\mathrm{ndf}=\chiPaper$.
It is computed with the pseudo-inverse over the 205 bins, which keeps all 205 eigenmodes (rank 205, condition
number $1.5\times10^{12}$), and ndf is that rank.
It indicates a correlated shape difference, not a normalization difference, spread over many eigenmodes of the
published covariance.
The source of the shape difference has not been identified.
Figure~\ref{fig:validation} shows the principal one-dimensional projections.

We quote the LightGBM result as the central value because its covariance uses the same estimator.
The distance depends on the estimator (\nsec{sec:results-tension}).
The learner was selected after the LightGBM and exact-split comparisons with the published result, so the
benchmark is a reproduction rather than an independent validation.
The purity-subtracted exact-split result differs from the central value by a median of \SI{1.03}{\percent} per bin
(maximum \SI{17.6}{\percent}), as detailed in \nsec{sec:method-config}.

\begin{figure*}[t]
\centering
\includegraphics[width=0.86\textwidth]{paper_fig1_validation}
\caption{Comparison with the published result~\cite{MINERvA:2021owq}. Top: $\pT$ and $\pPar$ projections of the
published result and of this analysis, each with its total $\pm1\sigma$ band, and of the MINERvA Tune~v1, with
ratios to the published result below. The band of this analysis contains the statistical, systematic-universe
(including the normalization), training-seed and run-period blocks of the covariance, Sec.~\ref{sec:uq2d}, and it
omits the learner dependence and $C_{\mathrm{model}}$. Here the two bands are correlated, so their overlap is
not a compatibility test. Bottom left: $(\mathrm{OmniFold}-\mathrm{published})/\sigma_{\mathrm{published}}$ per
reported bin, with cells drawn at equal size (grey: not reported). Bottom right: median relative uncertainties and
residual summary.}
\label{fig:validation}
\end{figure*}

\subsection{Two-dimensional covariance}
\label{sec:uq2d}

The two-dimensional covariance is the sum of four blocks, each from LightGBM unfoldings with the central value's
estimator: systematic universes (187 universes), statistical (300 bootstrap replicas), run period (three period
groups), and training seed (five seed settings, so that its own size is poorly determined).
The resulting relative uncertainty has a median of \SI{\uqMedian}{\percent} per bin and is
\SI{\uqIntegral}{\percent} on the integrated cross section, compared with \SI{\uqPaper}{\percent} and \SI{\uqPaperIntegral}{\percent} for the
publication.
The published flux uncertainty (median \SI{4.01}{\percent}) is larger than ours (\SI{\fluxBand}{\percent}), for a
reason the released inputs do not let us determine.
The covariance omits the approximation in the systematic-universe unfoldings for events that fail reconstruction,
the learner dependence of Sec.~\ref{sec:repro2d-pub}, and the dependence on the regularization, while the
dependence on the prior's generator shape is reported beside it (below).
The covariance has not been independently validated.
It provides no calibrated significance for generator comparisons.
The block sizes, the assumptions of the block sum, and the tests of the approximation are given in
Ref.~\cite{AnalysisNote2026}, Secs.~\ref*{N-sec:results-uq}, \ref*{N-sec:uq-combine}, and~\ref*{N-sec:uq-2d-scope}.

The statistical block, with a median statistical uncertainty of \SI{0.676}{\percent}, is a bootstrap of 300
replicas drawn at the central value's training seed, in which the background templates are also resampled.
The coverage of this block and of the total covariance has not been tested.
We make no coverage claim on independent simulated samples or simultaneously over all reported bins and
projections.

Data from different run periods are not mutually consistent within the other uncertainties.
One period group's unfolded $\pPar$ distribution, taken relative to its own simulation, differs from that of the
earliest periods by a step of about \SI{7}{\percent} near \SI{8}{GeV/c}.
The cause is not known, and these data cannot establish which period is correct.
The run-period block is built from double ratios among the three period groups (each group's unfolded result over
its own simulated truth, with no group chosen as the reference), includes this inconsistency in the combined result, and amounts to
\SI{2.0}{\percent} on the integrated cross section.
Because the block is built from the observed differences among the periods, its size is not validated and its own
consistency check cannot fail.

We report the dependence on the prior's generator shape as a separate covariance beside the four-block covariance,
$C_{\mathrm{model}}=\sum_k s_k s_k^{\mathsf{T}}$.
Here $s_k$ is the per-bin relative change of the result when the prior truth is reweighted by the ratio of the
$(\pT,\pPar)$ shape of generator~$k$ (GENIE with meson-exchange currents, NuWro, or GiBUU) to that of GENIE, with
its normalization kept.
Its median is \SI{\cmodelMedian}{\percent} per bin.
It is a measure of generator-shape model dependence, not a calibrated bound on the unfolding bias.
Added to the four-block covariance, it gives a total with a median of \SI{\qrCmodelMedian}{\percent} per bin.
The comparison with the published uncertainty uses the covariance without $C_{\mathrm{model}}$, because neither
has a component for the residual dependence on the prior's generator shape.
For the published analysis, this is because it chose its ten iterations in a fake-data study and assigns no uncertainty to the remaining
dependence~\cite{MINERvA:2021owq}.
The iteration counts differ (ten there, five here), so we make no statement about the size of the published
analysis's residual.

\subsection{Normalization and closure}
\label{sec:closure}

Of the joint-test estimator's results, only the three-observable marginal has been compared with the
lower-dimensional result it extends.
Its area-weighted total agrees with the two-dimensional result to \SI{0.141}{\percent}, a comparison across the
two estimators.
The normalization checks of the earlier purity-subtracted five-observable central value are described in
\nsec{sec:eavailw-axis}, and the closure checks of its estimator in three and four observables in
\nsec{sec:closure-map}.

Figure~\ref{fig:bkg-residual} shows that, with purity-based background subtraction in a deterministic LightGBM
estimator, end-to-end pseudo-experiments at nominal truth reproduce the input except in the highest-$W$ cells,
which are biased by up to \SI{4}{\percent} in simulation.
In the same estimator, unbinned negative-weight subtraction with a refinement of 100 trees and 8 leaves reduces it
to below \SI{0.3}{\percent} in those cells and at most \SI{0.4}{\percent} elsewhere in the projection, leaving up
to \SI{1}{\percent} in some joint cells.
On data, the two treatments differ by at most \SI{1.4}{\percent}.
Both kinds of closure study are described in \nsec{sec:closure-map}.

In two observables, a closure test with backgrounds (20 pseudo-experiments at nominal truth) finds the purity and
refined estimates more than one statistical standard deviation apart in 14 of the 205 cells.
The purity estimate lies up to 5.9 statistical standard deviations from the truth, and more than one total
standard deviation in one cell, while the refined one lies within 0.44 in every cell.
The two-dimensional result therefore uses the refined estimator (\napp{sec:negweight-footing}).

\begin{figure*}[t]
\centering
\includegraphics[width=0.84\textwidth]{paper_fig_bkg_residual}
\caption{Fractional closure residual, (unfolded $-$ truth)/truth, at nominal truth, for pseudo-experiments that
include backgrounds. Both from a deterministic LightGBM estimator. Circles: purity-based background subtraction
(9 pseudo-experiments). Squares: unbinned negative-weight subtraction with a refinement of 100 trees and 8 leaves
(60 pseudo-experiments). The pseudo-experiments are not paired. (a) The 42 $(\Eavail,W)$ projection cells, grouped
by $\Eavail$ bin. (b) The \jtNcells{} coarse joint cells, labelled by their index in the $3^{5}$ grid, then the
total integrated cross section. Error bars are the standard error of the mean.}
\label{fig:bkg-residual}
\end{figure*}

\subsection{Generator-sized truth changes}
\label{sec:departure}

Figure~\ref{fig:deformation} shows that the unfolding recovers only part of a change of the $\Eavail$ shape as
large as the difference between generators.
When the true shape is reweighted by the GiBUU-to-GENIE ratio, the $(\Eavail,W)$ cells unfolded with a
development variant of the joint-test estimator using a 100-tree, 8-leaf refinement deviate from it by several
percent in the largest cells and by up to \SI{74}{\percent} at high $W$, in the cell $0.8\le\Eavail<\SI{1.5}{GeV}$,
$W\ge\SI{3}{GeV}$.
These deviations occur with or without background, so background subtraction alone does not explain the residual.
We interpret the residual as regularization bias, because it persists over the tested iteration numbers and
classifier capacities, although no single cause is established.
For other fixed generator shapes tested with the joint-test estimator, it is \SIrange{16}{31}{\percent} in the
most affected $(\Eavail,W)$ cell and lies within the quoted total uncertainty in about half of those cells.
Those uncertainties come from an earlier purity-subtracted five-observable covariance that fails its specified
acceptance criteria (\nsec{sec:eavailw-adopted}).
For the sixteen-cell result, Sec.~\ref{sec:fd16} reports the fake-data tests of the estimator's bias.
The deformation studies, the iteration and capacity tests, the exact binned comparison and the
joint-test-estimator shapes are given in Ref.~\cite{AnalysisNote2026}, Secs.~\ref*{N-sec:val-negweight}
and~\ref*{N-sec:val-departure}.

\begin{figure*}[t]
\centering
\includegraphics[width=0.84\textwidth]{paper_fig_deformation}
\caption{Fractional recovery residual, (unfolded $-$ truth)/truth, per $(\Eavail,W)$ cell, where the truth is each
pseudo-experiment's deformed truth; all pseudo-experiments include backgrounds. (a) A development variant of the
joint-test estimator, whose refinement has 100 trees and 8 leaves, with the truth reweighted by the GiBUU/GENIE
$\Eavail$ shape ratio (60 pseudo-experiments), where the grey step is the injected change of the cell truth. At
nominal truth this variant's largest residual is \SI{0.39}{\percent} (the squares of Fig.~\ref{fig:bkg-residual}a).
(b) The joint-test estimator, for three shape changes not used while it was chosen (NuWro in $\Eavail$; GENIE + MEC
in $(\Eavail,W)$; NuWro in $(\pT,\pPar,\Eavail)$), 20 pseudo-experiments each, where the grey band is $\pm$ the
total uncertainty of the $(\Eavail,W)$ projection of the covariance of an earlier purity-subtracted
five-observable central value (\nsec{sec:eavailw-adopted}). At nominal truth (40 pseudo-experiments) its largest
residual is \SI{0.09}{\percent}. Neither seed set was held out from the estimator's development. $W$ bin 6 is
$W\ge\SI{3}{GeV}$. Error bars are the standard error of the mean.}
\label{fig:deformation}
\end{figure*}

A binned, signal-only and noise-free diagnostic on the \jtNcells{} joint cells and the truth changes used to
develop the estimator, with the detector response fixed at the simulation's, shows that exact binned iterative
unfolding with five iterations leaves about the same median residual as the boosted-tree estimator.
It covers only the residuals whose diagnosis is stable under the truth resolution (\SI{38}{\percent} of the
joint-cell residuals above \SI{2}{\percent}).
What unbinned information could recover, the coverage failure, and the validity of any interval lie outside its
scope.
Within that scope, the residual is consistent with the regularization of the inverse problem at this binning and
prior, rather than with the tree approximation of the classifiers (\nsec{sec:val-departure}).

\subsection{Coverage}
\label{sec:coverage}

Statistical intervals from a deterministic LightGBM estimator with purity-based subtraction, corrected for the
estimator's closure bias at nominal truth, miscover at the tested truths, most strongly when the true signal rate
is lowered by about \SI{10}{\percent}.
Under that change, 51 of the 153 reported cell and projection integrals have a 68\% coverage whose upper confidence
bound lies below the acceptance threshold of 0.66, with coverage as low as 0.06.
The end-to-end test unfolds 400 pseudo-experiments per truth.

A diagnosis attributes the miscoverage to a residual bias that grows with $W$ and survives the nominal-truth bias
correction.
The source of this residual bias has not been established.
Agreement in the marginal normalizations coexists with substantial residual bias under the tested deformations.
The interval construction and the diagnosis are described in \nsec{sec:coverage-5d}.

\subsection{Estimator stability}
\label{sec:cov5d}

The joint-test estimator is unstable at the level of rounding.
Rounding-level changes of its inputs move its results by up to 2.6 statistical standard deviations, and those of
the lower-capacity development variant by up to 2.5, so the estimator fails a data-stability requirement of at most 0.3 (median)
and 1.0 (maximum) statistical standard deviations.
This spread lies inside the data-bootstrap statistical uncertainty but limits how exactly the central values can be
reproduced.

\subsection{Fake-data tests of the sixteen-cell result}
\label{sec:fd16}

To compare the estimator's bias with the reported covariance when the truth differs from the simulation, we unfold
pseudo-data at \fvval{f0b-all-truths} truths and compare the mean unfolded result with each truth on the 16 cells
through $\chi^2$ with 16 degrees of freedom.
The truths are not held out: they are the simulation's prior; five hybrid truths, each a generator prediction on
the fine grid with the MINERvA Tune~v1 composition below it; NuWro shape changes in $\Eavail$ and in
$(\pT,\pPar,\Eavail)$; a GENIE + MEC shape change in $(\Eavail,W)$; a GiBUU-to-GENIE $\Eavail$ ratio; a noise-free
closure at the NuWro shape; and a synthetic tilt of the $q_3$ distribution at fixed $(\Eavail,W)$.
For a truth from a generator that contributes to $C_{\mathrm{model}}$, the covariance omits that generator's
component (the Cov. column of Table~\ref{tab:fd16}).
A truth is listed when $p<\fvval{trig-p}$ or when a cell has $|z|>\fvval{trig-z}$.
Table~\ref{tab:fd16} gives the results under the reported covariance and under $C_{\mathrm{total}}$ alone, and
\fvval{f0b-nungraded} further truths are ungraded and given in the note.

\begin{table*}[tb]
\caption{Fake-data tests on the sixteen cells: $\chi^2$ (16 degrees of freedom), $p$ and the cells with
$|z|>\fvval{trig-z}$, against $C_{\mathrm{total}}+C_{\mathrm{resp}}$ and against $C_{\mathrm{total}}$, each without
the truth's own generator component where the second column says so ($-$).}
\label{tab:fd16}
\begin{tabular}{@{}llrrlrrl@{}} \toprule
& & \multicolumn{3}{c}{$C_{\mathrm{total}}+C_{\mathrm{resp}}$} & \multicolumn{3}{c}{$C_{\mathrm{total}}$} \\ \cmidrule(lr){3-5}\cmidrule(l){6-8}
Truth & Cov. & $\chi^2$ & $p$ & cells & $\chi^2$ & $p$ & cells \\ \midrule
simulation prior & full & \fvval{f0b-nominal-chi2} & $\fvval{f0b-nominal-p}$ & \fvval{f0b-nominal-cells} & \fvval{f0-nominal-chi2} & $\fvval{f0-nominal-p}$ & \fvval{f0-nominal-cells} \\
Tune~v1, hybrid & full & \fvval{f0b-A7_mnvtune-chi2} & $\fvval{f0b-A7_mnvtune-p}$ & \fvval{f0b-A7_mnvtune-cells} & \fvval{f0-A7_mnvtune-chi2} & $\fvval{f0-A7_mnvtune-p}$ & \fvval{f0-A7_mnvtune-cells} \\
GENIE CV, hybrid & full & \fvval{f0b-A7_genie_cv-chi2} & $\fvval{f0b-A7_genie_cv-p}$ & \fvval{f0b-A7_genie_cv-cells} & \fvval{f0-A7_genie_cv-chi2} & $\fvval{f0-A7_genie_cv-p}$ & \fvval{f0-A7_genie_cv-cells} \\
GENIE + MEC, hybrid & $-$GENIE + MEC & \fvval{f0b-A7_genie_mec-chi2} & $\fvval{f0b-A7_genie_mec-p}$ & \fvval{f0b-A7_genie_mec-cells} & \fvval{f0-A7_genie_mec-chi2} & $\fvval{f0-A7_genie_mec-p}$ & \fvval{f0-A7_genie_mec-cells} \\
GENIE + MEC, $(\Eavail,W)$ shape & $-$GENIE + MEC & \fvval{f0b-W2-chi2} & $\fvval{f0b-W2-p}$ & \fvval{f0b-W2-cells} & \fvval{f0-W2-chi2} & $\fvval{f0-W2-p}$ & \fvval{f0-W2-cells} \\
NuWro, hybrid & $-$NuWro & \fvval{f0b-A7_nuwro-chi2} & $\fvval{f0b-A7_nuwro-p}$ & \fvval{f0b-A7_nuwro-cells} & \fvval{f0-A7_nuwro-chi2} & $\fvval{f0-A7_nuwro-p}$ & \fvval{f0-A7_nuwro-cells} \\
NuWro, $\Eavail$ shape & $-$NuWro & \fvval{f0b-W1-chi2} & $\fvval{f0b-W1-p}$ & \fvval{f0b-W1-cells} & \fvval{f0-W1-chi2} & $\fvval{f0-W1-p}$ & \fvval{f0-W1-cells} \\
NuWro, $(\pT,\pPar,\Eavail)$ shape & $-$NuWro & \fvval{f0b-W3-chi2} & $\fvval{f0b-W3-p}$ & \fvval{f0b-W3-cells} & \fvval{f0-W3-chi2} & $\fvval{f0-W3-p}$ & \fvval{f0-W3-cells} \\
NuWro, closure & $-$NuWro & \fvval{f0b-closure_nuwro-chi2} & $\fvval{f0b-closure_nuwro-p}$ & \fvval{f0b-closure_nuwro-cells} & \fvval{f0-closure_nuwro-chi2} & $\fvval{f0-closure_nuwro-p}$ & \fvval{f0-closure_nuwro-cells} \\
GiBUU, $\Eavail$ shape & $-$GiBUU & \fvval{f0b-stage0c_gibuu-chi2} & $\fvval{f0b-stage0c_gibuu-p}$ & \fvval{f0b-stage0c_gibuu-cells} & \fvval{f0-stage0c_gibuu-chi2} & $\fvval{f0-stage0c_gibuu-p}$ & \fvval{f0-stage0c_gibuu-cells} \\
GiBUU, hybrid & $-$GiBUU & \fvval{f0b-A7_gibuu-chi2} & $\fvval{f0b-A7_gibuu-p}$ & \fvval{f0b-A7_gibuu-cells} & \fvval{f0-A7_gibuu-chi2} & $\fvval{f0-A7_gibuu-p}$ & \fvval{f0-A7_gibuu-cells} \\
synthetic $q_3$ tilt & full & \fvval{f0b-q3-chi2} & $\fvval{f0b-q3-p}$ & \fvval{f0b-q3-cells} & \fvval{f0-q3-chi2} & $\fvval{f0-q3-p}$ & \fvval{f0-q3-cells} \\ \bottomrule
\end{tabular}
\end{table*}

Under the reported covariance, without the GiBUU component for the two GiBUU truths, the listed truths are the
GiBUU $\Eavail$ shape (cell R06, $z=\fvval{f0b-stage0c_gibuu-zmax}$), the GiBUU hybrid truth
($p=\fvval{f0b-A7_gibuu-p}$; cells \fvval{f0b-A7_gibuu-cells}), and the synthetic $q_3$ tilt
($p=\fvval{f0b-q3-p}$).
Thus the reported uncertainty does not cover the fake-data bias at these GiBUU-like truths or at the $q_3$ tilt.
At the other \fvval{f0b-nunlisted} truths the bias passes both criteria, with no cell beyond
$\fvval{trig-z}\sigma$ and $p\geq\fvval{trig-p}$.

Two GENIE~v3 tunes, \fvval{f1-tune-g18} and \fvval{f1-tune-ar23}, were generated for this test and held out of the
analysis.
Each reweights the simulated truth using its shape ratio to GENIE CV, constructed as for the $C_{\mathrm{model}}$
components.
The simulated detector response is not probed.
For each tune, \fvval{f1-n} Poisson pseudo-experiments and a noise-free closure are compared with the truth using
the full reported covariance.
Both tunes satisfy the bias criteria: the pseudo-experiment means give $\chi^2=\fvval{f1-g18-chi2}$
($p=\fvval{f1-g18-p}$) and \fvval{f1-ar23-chi2} ($p=\fvval{f1-ar23-p}$), respectively, with $|z|\leq\fvval{f1-zmax-ceil}$ in every
cell of the means and closures.
In cell R11 at \fvval{f1-tune-g18}, the bias lies below the reported uncertainty only with $C_{\mathrm{model}}$
included.
The tunes share a quasielastic cross-section model, and \fvval{f1-tune-g18} and the GENIE + MEC component of
$C_{\mathrm{model}}$ also use Valencia-family MEC models.
These tests provide partial evidence and are not independent.
Omitting that $C_{\mathrm{model}}$ component gives $\chi^2=\fvval{f1-g18-sens-chi2}$ ($p=\fvval{f1-g18-sens-p}$) and
\fvval{f1-ar23-sens-chi2} ($p=\fvval{f1-ar23-sens-p}$), respectively, as a descriptive comparison.
The tests assess bias, not coverage: only event counts fluctuate, with systematic parameters fixed, so interval
widths cannot be tested.
